%% ============================================================
%%  SOLUCIÓN — Ejercicio 5: Paper completo (INTEGRADOR)
%%  Clase 2 — Al final de la Clase 2
%%  Compilación: pdflatex → biber → pdflatex → pdflatex
%%  (O en TeXstudio: F5 → F11 → F5 → F5)
%% ============================================================

\documentclass[12pt, a4paper]{article}

% --- Codificación ---
\usepackage[utf8]{inputenc}
\usepackage[T1]{fontenc}
\usepackage[spanish]{babel}
\usepackage{lmodern}

% --- Geometría ---
\usepackage[top=2.5cm, bottom=2.5cm, left=3cm, right=3cm]{geometry}

% --- Matemáticas y teoremas ---
\usepackage{amsmath, amssymb, amsthm}
\usepackage{booktabs}
\usepackage{graphicx}
\usepackage[colorlinks=true, linkcolor=black, citecolor=blue]{hyperref}

% --- MACROS ECONOMÉTRICAS ---
\newcommand{\E}{\mathbb{E}}
\newcommand{\Var}{\mathrm{Var}}
\newcommand{\Cov}{\mathrm{Cov}}
\newcommand{\plim}{\operatorname{plim}}

% --- ENTORNOS MATEMÁTICOS ---
\newtheorem{theorem}{Teorema}[section]
\newtheorem{proposition}[theorem]{Proposición}
\newtheorem{lemma}[theorem]{Lema}

\theoremstyle{definition}
\newtheorem{assumption}{Supuesto}

\theoremstyle{remark}
\newtheorem{remark}{Observación}

% --- BIBLIOGRAFÍA ---
\usepackage[style=authoryear, backend=biber, maxnames=3]{biblatex}
\addbibresource{referencias.bib}

% --- METADATOS ---
\title{Salarios y educación en Paraguay:\\Una aplicación de la ecuación de Mincer}
\author{Tu Nombre}
\date{\today}

% ================================================================
\begin{document}
% ================================================================

\maketitle
\thispagestyle{empty}

% --- Abstract ---
\begin{abstract}
\noindent
Este trabajo estima los retornos a la educación para Paraguay usando datos 
de la Encuesta Permanente de Hogares 2023 y la ecuación de Mincer. Estimamos 
que un año adicional de educación aumenta el salario en aproximadamente 
6.9\%, controlando por experiencia laboral y características demográficas. 
Documentamos una importante brecha de género: las mujeres ganan 18.3\% menos 
que los hombres con igual educación y experiencia. Estos resultados sugieren 
que la inversión en educación es económicamente rentable, pero que persisten 
disparidades en el mercado laboral que requieren intervención de política pública.

\vspace{0.5em}
\noindent \textbf{Palabras clave:} capital humano, retornos a la educación, 
brecha de género, mercado laboral, Paraguay.

\noindent \textbf{Clasificación JEL:} J24, J31, I20.
\end{abstract}

\newpage
\tableofcontents
\newpage

% ================================================================
\section{Introducción}
% ================================================================

Desde Becker (1964), sabemos que la educación es una forma de inversión en 
capital humano que aumenta la productividad y, por tanto, los salarios. La 
\textcite{mincer1974} proporciona un marco flexible para estimar los retornos 
a la educación en datos de corte transversal.

En Paraguay, poco se conoce sobre los retornos a la educación o sobre cómo 
varían estos retornos según características individuales. Este trabajo llena 
ese vacío estimando ecuaciones de ingresos para una muestra representativa 
de trabajadores.

\subsection{Contribuciones}

Este trabajo contribuye de tres formas:

\begin{enumerate}
  \item Proporciona estimaciones de los retornos a la educación específicas 
        para Paraguay, contexto que carece de evidencia reciente.
  \item Documenta brechas de género en salarios y analiza su origen.
  \item Usa métodos econométricos estándar (OLS con errores robustos) que 
        sirven como base para investigación futura con técnicas más sofisticadas.
\end{enumerate}

% ================================================================
\section{Marco teórico y empírico}
% ================================================================

\subsection{El modelo de capital humano}

El modelo de capital humano postula que los individuos invierten en educación 
porque esperan retornos futuros en la forma de salarios más altos. La ecuación 
de Mincer, en su forma estándar, es:

\begin{equation}
  \ln(w_i) = \alpha + \beta_1 \, \text{educ}_i + \beta_2 \, \text{exp}_i 
             + \beta_3 \, \text{exp}_i^2 + \gamma \, \text{mujer}_i + \varepsilon_i
  \label{eq:mincer}
\end{equation}

donde $w_i$ es el salario mensual y $\varepsilon_i$ es el término de error.

\subsection{Supuestos y estimación}

\begin{assumption}
  \label{ass:linearity}
  La relación entre log-salarios e inputs educativos es lineal en los parámetros.
\end{assumption}

\begin{assumption}
  \label{ass:ceteris}
  \label{ass:exogeneity}
  Exogeneidad: $\E[\varepsilon_i | \text{educ}_i, \text{exp}_i, \text{mujer}_i] = 0$. 
  Es decir, educación no está correlacionada con factores no observados que 
  afecten salarios (habilidad, motivación).
\end{assumption}

\begin{proposition}
  \label{prop:ols_consistency}
  Bajo los Supuestos \ref{ass:linearity}--\ref{ass:exogeneity}, el estimador 
  OLS de la ecuación \eqref{eq:mincer} es consistente.
\end{proposition}

\begin{remark}
  El Supuesto \ref{ass:exogeneity} es fuerte y probablemente se viola en la 
  práctica: habilidad no observada está correlacionada tanto con educación 
  (personas más capaces estudian más) como con salarios. Esto genera sesgo 
  de variable omitida. Métodos como variables instrumentales podrían abordar 
  este problema.
\end{remark}

% ================================================================
\section{Datos}
% ================================================================

Los datos provienen de la Encuesta Permanente de Hogares (EPH) 2023, realizada 
por el Instituto Nacional de Estadística. La EPH es una encuesta de hogares 
de propósito general que cubre información sobre empleo, educación, y 
características demográficas.

Para este análisis, seleccionamos individuos ocupados de 18 a 65 años con 
información completa en todas las variables necesarias. La muestra final 
contiene 1.200 observaciones.

\begin{table}[htbp]
  \centering
  \caption{Estadísticas descriptivas}
  \label{tab:desc}
  \begin{tabular}{lrrrr}
    \toprule
    Variable                    & Media  & D.E.   & Mín.  & Máx.  \\
    \midrule
    Log(salario mensual)        & 14.52  & 0.53   & 13.22 & 16.10 \\
    Años de educación           & 12.3   & 3.4    & 0     & 20    \\
    Experiencia potencial       & 8.7    & 5.1    & 0     & 40    \\
    Mujer (=1)                  & 0.48   & 0.50   & 0     & 1     \\
    \midrule
    Observaciones               & \multicolumn{4}{c}{1.200} \\
    \bottomrule
  \end{tabular}
  \begin{minipage}{\linewidth}
    \footnotesize \textit{Notas:} La experiencia potencial se calcula como 
    edad $-$ años de educación $-$ 6. \textit{Fuente:} EPH 2023, INE.
  \end{minipage}
\end{table}

% ================================================================
\section{Resultados}
% ================================================================

La Tabla \ref{tab:regresion} presenta nuestras estimaciones. La columna (1) 
estima el retorno a la educación sin otros controles. La columna (2) agrega 
experiencia y su cuadrado. La columna (3) incluye un indicador de género.

\begin{table}[htbp]
  \centering
  \caption{Ecuación de Mincer: estimaciones OLS}
  \label{tab:regresion}
  \begin{tabular}{lccc}
    \toprule
                          & (1)        & (2)        & (3)        \\
                          & \multicolumn{3}{c}{\textit{Variable dep.: } $\ln(\text{salario})$} \\
    \cmidrule(lr){2-4}
    \midrule
    Educación             & 0.082***   & 0.071***   & 0.069***   \\
                          & (0.008)    & (0.009)    & (0.010)    \\[4pt]
    Experiencia           &            & 0.034**    & 0.031**    \\
                          &            & (0.012)    & (0.013)    \\[4pt]
    Experiencia$^2$       &            & $-$0.001** & $-$0.001*  \\
                          &            & (0.000)    & (0.000)    \\[4pt]
    Mujer (=1)            &            &            & $-$0.183***\\
                          &            &            & (0.042)    \\[4pt]
    Constante             & 7.241***   & 7.018***   & 7.312***   \\
                          & (0.098)    & (0.115)    & (0.122)    \\
    \midrule
    Observaciones         & 1.200      & 1.200      & 1.200      \\
    $R^2$ ajustado        & 0.241      & 0.318      & 0.374      \\
    \bottomrule
  \end{tabular}
  \begin{minipage}{\linewidth}
    \footnotesize \textit{Notas:} Errores estándar robustos a 
    heterocedasticidad en paréntesis. *** $p<0.01$, ** $p<0.05$, 
    * $p<0.1$. \textit{Fuente:} EPH 2023, elaboración propia.
  \end{minipage}
\end{table}

\subsection{Interpretación}

El coeficiente de educación en la especificación (3) es 0.069, lo que implica 
que un año adicional de educación aumenta el salario en aproximadamente 
6.9\%. Este retorno es estadísticamente significativo y económicamente 
importante.

Analizando el sistema de ecuaciones con experiencia:
\begin{align}
  \ln(w_i) &= 7.312 + 0.069 \cdot \text{educ}_i + 0.031 \cdot \text{exp}_i 
             - 0.001 \cdot \text{exp}_i^2 - 0.183 \cdot \text{mujer}_i + \varepsilon_i
\end{align}

El ciclo de vida de ingresos tiene forma de U invertida: los ingresos suben 
con experiencia pero a tasa decreciente. El retorno máximo a la experiencia 
ocurre a aproximadamente $\text{exp}^* = 0.031 / (2 \times 0.001) = 15.5$ años 
de experiencia.

La brecha de género es substantiva: ser mujer reduce el log-salario en 0.183, 
o equivalentemente, reduce el salario en $(e^{0.183} - 1) \approx 20\%$ en 
términos de nivel (no logarítmico).

% ================================================================
\section{Conclusiones}
% ================================================================

Este trabajo estima que los retornos a la educación en Paraguay son de 
aproximadamente 7\% por año de escolaridad. Este retorno es comparable a 
estimaciones internacionales y sugiere que la inversión en educación es 
económicamente rentable.

Sin embargo, documentamos una importante brecha de género que persiste incluso 
después de controlar por educación y experiencia. Esta brecha podría reflejar:
\begin{enumerate}
  \item Discriminación en el mercado laboral.
  \item Segregación ocupacional (mujeres en ocupaciones de menor remuneración).
  \item Diferencias en negociación de salarios.
  \item Otras fuentes de heterogeneidad no observada.
\end{enumerate}

Investigaciones futuras deberían explorar los mecanismos detrás de esta brecha 
y estimar modelos con mayor sofisticación econométrica.

\newpage
\printbibliography[title={Referencias}]

\end{document}
